\chapter{Introdução}
\label{capitulo:introducao}

No capítulo da introdução é feita a definição inicial do tema da pesquisa ou trabalho a ser desenvolvido.
Neste momento, o autor do trabalho deve ser claro e sua explicação deve ser simples evitando o máximo possível de termos técnicos.

Pode-se dizer que uma boa introdução ou até mesmo um bom texto é aquele que qualquer leitor, independente do conhecimento prévio na área, consegue entender bem.

Este texto introdutório apresenta o \textbf{tema} e o \textbf{problema de pesquisa}.
No entanto, não deve ser um texto muito longo \cite{RaulSidneiWazlawick2021}.
Recomendo que este trecho tenha no máximo 2 páginas.

\section{Objetivo Geral}

Em um parágrafo apresentar o objetivo do trabalho. Por exemplo: ``Este trabalho tem como objetivo desenvolver um sistema capaz de resolver o problema XYZ com um tempo polinomial''.

\section{Objetivos específicos}

Para concluir o objetivo geral foram definidos os seguintes objetivos específicos: 

\begin{itemize}
	\item Determinar uma arquitetura que permita comunicação entre as partes do sistema;

    \item Desenvolver um sistema administrativo que permita cadastrar dados e emitir relatórios de acompanhamento;

    \item Avaliar a usabilidade do sistema através do \textit{feedback} do usuário e propor melhorias nas interfaces gráficas.
\end{itemize}

\section{Justificativa}
\label{secao:justificativa}

É nesta seção que você deverá justificar a importância do trabalho.
Deve-se deixar claro o motivo deste trabalho ser relevante e enfatizar o que ele difere dos demais trabalhos existentes que serão apresentados no capítulo de revisão bibliográfica.

\section{Contribuições realizadas}

Esta seção entra com pontos muito bem objetivos sobre o que foi avançado na área de pesquisa do tema do trabalho.
Note que para trabalhos que envolvam apenas o desenvolvimento de um \textit{software} sem a apresentação de grandes melhorias ou novas propostas no que já existe.
Este pode ser um ponto crucial para o aceite ou aprovação da banca.

A seção \ref{secao:justificativa} pode te apoiar na escrita das contribuições realizadas, uma vez que, se sua crítica é sobre o tempo de resposta de um determinado algoritmo, espera-se que o trabalho em questão tenha melhora este tempo.

\section{Organização do trabalho}

Este capítulo expôs o problema de pesquisa, os motivos, os objetivos e as contribuições deste trabalho.
O capítulo \ref{capitulo:revisao_bibliografica} apresenta os principais trabalhos relacionados ao tema assim como as principais ferramentas de \textit{frameworks} utilizados para a solução do problema.
O capítulo \ref{capitulo:materiais_metodos} apresenta os materiais e métodos utilizados como abordagem desta pesquisa.
O capítulo \ref{capitulo:projeto} é apresentado o projeto de \textit{software} proposto neste trabalho, assim são apresentados diagramas relacionados a arquitetura, banco de dados, estrutura de fluxo de comunicação entre as entidades que compõem o sistema.
O capítulo \ref{capitulo:resultados_discussao} são apresentados os resultados deste trabalho e faz-se uma discussão sobre os motivos destes serem como são de suas consequências na área de pesquisa.
Por fim, o capítulo \ref{capitulo:consideracoes_finais} apresenta as considerações finais e propostas de trabalhos futuros relacionados ao tema.
